\begin{proposition}[DP27: parity and endpoint constraints of the existing chain]\label{DP27}
In \path{Formalism/LocalRingCardinality/Params.lean}, \path{Ranks} requires an even lift count $M\ge2$. With its $k,t$ fixed, its endpoint ranks satisfy
\[
 N(M)-N(L)=(M-L)(t+1)>0\qquad(0\le L<M).
\]
Thus a replacement by $M-1$ violates the same parameter specification, and literal truncation cannot preserve the old hermitian endpoint by an isometry. The auxiliary signature at $1\le L<M$ is $(k+L-1,M-L)$, whereas at $M$ it is $(k+M-1,0)$. The column spaces $X_j$ themselves are definite at every stage. A truncated column does not, by definition, supply the period identification at the old endpoint.
\end{proposition}
\begin{proof}
The parity statement is the field \path{Ranks.m_even}; the difference between two allowed even counts is even. The definitions are $N(j)=k+j(t+1)$ and $\operatorname{sig}W(j)=(k+j-1,M-j)$. Since $t+1>0$, the displayed rank difference is positive. The source field $k\ge M+3$ and $1\le L<M$ make both entries of the intermediate auxiliary signature positive. These numerical ranks and signatures are imposed on the actual spaces by \path{DataSpaces.lean}, \path{Spaces.X_rank}, \path{W_sigval}; its separate \path{X_def} requires all column spaces to be definite. Rank is invariant under an isometry, proving the endpoint obstruction.

Finally \path{DataKernels.lean}, \path{ColumnChain.col}, is \path{chainOpM} at exactly $M$, with output \path{C.cls M}\,$\to$\,\path{C.Y M}\,$\to\mathbf C$. The field \path{PeriodMap.evz} consumes that terminal space. Truncation gives stage $L$ instead. A numerical lower bound on values supplies none of the new representation or period-comparison data. This is a limitation of literal truncation, not a prohibition on separately constructing admissible data with a smaller even count.
\end{proof}
